% Generated draft: requires mathematical and source review.
\documentclass[11pt]{article}
\usepackage[margin=25mm]{geometry}
\usepackage{fontspec}
\setmainfont{Latin Modern Roman}
\usepackage{amsmath,amssymb,mathtools,unicode-math}
\setmathfont{Latin Modern Math}
\usepackage{enumitem,needspace,xurl,hyperref}
\setlength{\parindent}{0pt}
\setlength{\parskip}{5pt}
\setlength{\emergencystretch}{3em}
\newcommand{\N}{\mathbb{N}}
\newcommand{\Z}{\mathbb{Z}}
\newcommand{\Q}{\mathbb{Q}}
\newcommand{\R}{\mathbb{R}}
\newcommand{\C}{\mathbb{C}}
\newcommand{\F}{\mathbb{F}}
\newcommand{\A}{\mathbb{A}}
\newcommand{\PP}{\mathbb{P}}
\newcommand{\E}{\mathbb{E}}
\newcommand{\eps}{\varepsilon}
\newcommand{\dd}{\,\mathrm{d}}
\begin{document}
\section*{UnsolvedMath ID: 800031}
\textbf{Source problem code:} catalogue entry
\subsection{Definitions and mathematical statement}
A semigroup is a set equipped with an associative binary operation. Let \(\Phi\) be an arbitrary finite class of finite semigroups, and let
\[
\mathbb K\in\{\mathbb H\mathbb P,\mathbb H\mathbb P\mathbb S,\mathbb S\mathbb H\mathbb P,\mathbb S\mathbb H\mathbb P\mathbb S\}.
\]
Here the symbols are the Birkhoff operators used in the source. Conventionally, \(\mathbb H\) denotes taking homomorphic images, \(\mathbb S\) taking subsemigroups, and \(\mathbb P\) taking direct products; the supplied excerpt does not specify the permitted product index sets or the treatment of empty products.

The question attributed to S.\,I. Kublanovsky is: for every such \(\mathbb K\) and \(\Phi\), is the class
\[
\operatorname{Cor}\bigl(\mathbb K(\Phi)\bigr)
\]
decidable?

The supplied record does not define \(\operatorname{Cor}\), so this notation is retained without an invented interpretation. Consequently, neither the objects whose membership is to be decided nor their effective representation can be specified from the supplied evidence. If the intended membership inputs are finite semigroups presented by multiplication tables, decidability would mean that an algorithm halts on every such input and correctly determines membership in the displayed class; this is a conditional interpretation, not a hypothesis supplied by the source. The question is stated separately for each \(\mathbb K\) and \(\Phi\); whether a single algorithm must also accept \(\mathbb K\) and \(\Phi\) as input is not specified.

The source additionally states that the answer is positive for all other Birkhoff operators, referring to \([\mathrm{Kubl6}]\). The excerpt supplies neither the definition of that full family of operators nor the bibliographic entry for this reference. This contextual assertion is retained as a source claim.
\subsection{Short English statement}
For each of four specified Birkhoff operators and every finite class of finite semigroups, ask whether membership in the resulting \(\operatorname{Cor}\) class is decidable. The operation \(\operatorname{Cor}\) is undefined in the supplied excerpt, so the intended decision problem requires clarification. The source attributes the question to S.\,I. Kublanovsky and reports positive answers for the other Birkhoff operators.
\subsection{Sources}
\begin{itemize}
\item Local UnsolvedMath registry, permanent ID \(800031\), exact source problem code \texttt{catalogue entry}, title: \emph{The Lyapin's Notebook --- Problem 33}. The supplied record labels the problem open; this status is not independently verified.
\item Bershadsky, Kublanovsky, and Mashevitzky, \emph{The Lyapin's Notebook: Unsolved Problems in Semigroup Theory} (2026), Problem 33. Supplied source URL: \url{https://arxiv.org/abs/2604.04763}. The statement specifically attributes the question to S.\,I. Kublanovsky. The registry lists the three compilers in its proposer field and records 2026 as the proposed year; these metadata do not establish the original proposal date.
\item The supplied provenance records access on 2026-10-16 and extraction by \texttt{arxiv\_latex}. This is supplied provenance, not a claim of independent source access or verification.
\item Reference label \([\mathrm{Kubl6}]\), cited in the source for the assertion concerning all other Birkhoff operators. No bibliographic details or separate URL for this reference are supplied.
\end{itemize}
\end{document}
